\SourceLeaf{387}%
Nummer (1810) hörte sie ganz auf. Von dem General \textsc{Guilleminot}, welcher im Jahr 1822 die Direction des Depot übernahm, wurde zuerst die Idee einer regelmässigen Fortsetzung der Zeitschrift gefasst, und von dessen interimistischem Stellvertreter, dem General \textsc{Delachasse de Verigny} zur Reife gebracht. Man beschloss zugleich einige Abänderung in der Form eintreten, und in dem für die Fortsetzung gewählten Quartformat auch die sieben ältern Stücke. wovon die Exemplare vergriffen waren, von neuem abdrucken zu lassen. Zu diesem neuen Abdruck der früheren Stücke sind die beiden ersten Bände der neuen Ausgabe bestimmt, wovon der erste nebst drei Bänden der Fortsetzung vor uns liegt; der zweite soll nächstens nachgeliefert werden.

Das \textit{Dépôt de la guerre} wurde zuerst 1688 unter \textsc{Ludwig} XIV. Regierung durch den Minister \textsc{Louvois} gestiftet: es war jedoch Anfangs nur ein Archiv, in welchem die gesammte officielle Armee-Correspondenz hinterlegt wurde. Zu seiner gegenwärtigen Einrichtung ist es erst nach und nach durch Erweiterung seines Umfangs und Consolidirung seiner Organisation gelangt, und seinen eigenthümlich grossartigen Character hat es erst erhalten, seitdem es der Mittelpunkt geworden ist, wo sich alle Früchte der Arbeiten eines selbstständigen Corps, der Ingenieurs-Geographen, vereinigen. Einen Begriff von dem Umfange der Thätigkeit des Instituts gibt der Umstand, dass die jährlichen Kosten im Jahr 1801 auf 110000 Franken angeschlagen wurden, worin die Gehalte des Personals nicht mit begriffen waren; letztere betrugen 1793 die Summe von 231100 Franken.

An eine Zeitschrift, welche aus einer so überschwenglich reichen Quelle schöpfen kann, darf man grosse Ansprüche machen, und diese werden um so vollkommner befriedigt werden, je mehr die Herausgeber ihr Hauptaugenmerk auf die Bekanntmachung wichtiger, dem Publicum bisher verschlossener Materialien richten, und dasjenige, wodurch die Wissenschaft nicht weiter gebracht wird, ausschliessen werden.

Von diesem Ideal finden wir die ältern Artikel viel weiter entfernt, als die neuere Fortsetzung. In der That sind die Artikel des ersten Bandes, wenn wir einen Aufsatz über die Hydrographie eines Theils von Frankreich und eine Notiz über die Geschichte des \textit{Dépôt de la guerre} ausnehmen, von der Art, dass sie eben so gut hätten geschrieben werden können, wenn auch das \textit{Dépôt} gar nicht vorhanden gewesen wäre, und ohne das Interesse zu leugnen, welches mehrere Aufsätze vor dreissig Jahren haben konnten und zum Theil noch jetzt haben, kann
